\documentclass[10pt,a4paper]{amsart}
\usepackage[margin=19mm]{geometry}
\usepackage{amsmath,amssymb,mathtools}
\usepackage[scr=rsfs,cal=euler]{mathalfa}
\usepackage{xfrac}
\usepackage[unicode,hidelinks,bookmarks=false]{hyperref}
\newcommand{\ie}{i.e.\ }
\newcommand{\cf}{cf.\ }
\newcommand{\resp}{respectively\ }
\newcommand{\Subsection}{\subsection}
\providecommand{\nobreakdash}{\nobreak\hbox{-}}
\setlength{\emergencystretch}{2em}
\multlinegap0pt
\usepackage{xcolor}
\usepackage{multirow}
\usepackage{booktabs}
\usepackage{tikz}
\usepackage{subfig}
\usepackage{amssymb}
\usepackage{algorithm}\usepackage{algpseudocode}
\newtheorem{lem}{Lemma}
\title{The Sequence Reconstruction of Permutations under Hamming Metric with Small Errors}
\author{A. Abdollahi, J. Bagherian, H. Eskandari, F. Jafari, M. Khatami, F. Parvaresh, R. Sobhani}
\date{Source: arXiv:2601.02844v2 (CC BY 4.0)}
\begin{document}
\maketitle

\noindent\emph{Source and licence.} The Sequence Reconstruction of Permutations under Hamming Metric with Small Errors. Version arXiv:2601.02844v2 (CC BY 4.0), distributed by arXiv under the Creative Commons Attribution 4.0 International licence, http://creativecommons.org/licenses/by/4.0/, as recorded in the arXiv licence metadata for this exact version and shown on the abstract page https://arxiv.org/abs/2601.02844v2. Full licence notice: \texttt{licenses/arxiv-2601.02844-CC-BY-4.0.txt}.\par\smallskip

\noindent\textit{Editorial scope.} Counted unit: the article's lem lem4 (source0.tex lines 215--218). The article's complete original proof of it is delivered with it (source0.tex lines 219--226, one \texttt{proof} environment of 8 lines). No mathematics is written, repaired or re-derived: the statement and the proof are the article's own lines, in the article's own notation, and no retained line is substituted or re-flowed.  No further source-local context is needed: every cross-reference of every retained line resolves inside the two delivered spans.  Nothing was omitted from any retained span, so the package carries no omission record.  No retained line contains a citation, so the wrapper carries no bibliography and no bibliography entry.  No retained line contains a figure environment, an image-inclusion command or drawing code, so no image is delivered and the package declares no asset dependency.  Editorial formatting only, in addition: an AMS \texttt{amsart} preamble that restates the article's own declarations and macros used by the retained lines, namely the environments \texttt{lem} on separate counters, together with the standard packages those lines need.  The retained environments print in this document as Lemma 1 in order of appearance, whereas the article prints the same items with the numbering of its own full document.  The complete native source of the article as distributed in the arXiv e-print for this exact version is bundled unchanged as \texttt{sources/192185/source0.tex}.
\begin{lem}\label{lem4}
	For any two permutations $\sigma$ and $\pi$ in $S_n$, $$C_{S_n}(\sigma, \pi)=C_{S_{M(\sigma)\cup M(\pi)}}(\sigma,\pi) \times 
		S_{[n] \setminus \big( M(\sigma) \cup M(\pi) \big)}.$$
\end{lem}

\begin{proof}
It suffices to prove that for each $\tau \in C_{S_n}(\sigma, \pi)$ and each $i \in M(\sigma) \cup M(\pi)$, we have $i^\tau \in M(\sigma) \cup M(\pi)$. 
Suppose  that $i^\tau = j \in [n] \setminus \big( M(\sigma) \cup M(\pi) \big)$. Without loss of generality, we may assume that  $i\in M(\sigma)$.
Since $\tau \in C_{S_n}(\sigma)$, we have $i^{\sigma\tau} = i^{\tau\sigma}$. 
Suppose that $i^\sigma = k \in M(\sigma) \setminus \{i\}$. 
Then we would have $i^{\tau} = k^\tau = j$, which is a contradiction. 
This completes the proof.
\end{proof}

\end{document}
